% Options for packages loaded elsewhere
\PassOptionsToPackage{unicode}{hyperref}
\PassOptionsToPackage{hyphens}{url}
\PassOptionsToPackage{dvipsnames,svgnames,x11names}{xcolor}
%
\documentclass[
  10pt,
  a4paper,
]{article}
\usepackage{amsmath,amssymb}
\usepackage{setspace}
\usepackage{iftex}
\ifPDFTeX
  \usepackage[T1]{fontenc}
  \usepackage[utf8]{inputenc}
  \usepackage{textcomp} % provide euro and other symbols
\else % if luatex or xetex
  \usepackage{unicode-math} % this also loads fontspec
  \defaultfontfeatures{Scale=MatchLowercase}
  \defaultfontfeatures[\rmfamily]{Ligatures=TeX,Scale=1}
\fi
\usepackage{lmodern}
\ifPDFTeX\else
  % xetex/luatex font selection
  \setmainfont[]{Liberation Serif}
  \setsansfont[]{Liberation Sans}
  \setmonofont[]{DejaVu Sans Mono}
\fi
% Use upquote if available, for straight quotes in verbatim environments
\IfFileExists{upquote.sty}{\usepackage{upquote}}{}
\IfFileExists{microtype.sty}{% use microtype if available
  \usepackage[]{microtype}
  \UseMicrotypeSet[protrusion]{basicmath} % disable protrusion for tt fonts
}{}
\makeatletter
\@ifundefined{KOMAClassName}{% if non-KOMA class
  \IfFileExists{parskip.sty}{%
    \usepackage{parskip}
  }{% else
    \setlength{\parindent}{0pt}
    \setlength{\parskip}{6pt plus 2pt minus 1pt}}
}{% if KOMA class
  \KOMAoptions{parskip=half}}
\makeatother
\usepackage{xcolor}
\usepackage[margin=22mm]{geometry}
\setlength{\emergencystretch}{3em} % prevent overfull lines
\providecommand{\tightlist}{%
  \setlength{\itemsep}{0pt}\setlength{\parskip}{0pt}}
\setcounter{secnumdepth}{5}
\ifLuaTeX
\usepackage[bidi=basic]{babel}
\else
\usepackage[bidi=default]{babel}
\fi
\babelprovide[main,import]{korean}
\ifPDFTeX
\else
\babelfont{rm}[]{Liberation Serif}
\fi
% get rid of language-specific shorthands (see #6817):
\let\LanguageShortHands\languageshorthands
\def\languageshorthands#1{}
\usepackage{xeCJK}
\setCJKmainfont{Noto Serif CJK KR}
\setCJKsansfont{Noto Sans CJK KR}
\setCJKmonofont{Noto Sans Mono CJK KR}
\xeCJKsetup{CJKspace=true}
\usepackage{xurl}
\usepackage{fvextra}
\usepackage{fancyhdr}
\AtBeginDocument{\DefineVerbatimEnvironment{verbatim}{Verbatim}{breaklines=true,breakanywhere=true,fontsize=\small}}
\setlength{\emergencystretch}{4em}
\setlength{\headheight}{15pt}
\pagestyle{fancy}\fancyhf{}
\fancyhead[L]{\small 중성자별 구조와 비방사 진동}
\fancyfoot[C]{\thepage}
\renewcommand{\headrulewidth}{0.3pt}
\setcounter{tocdepth}{2}
\ifLuaTeX
  \usepackage{selnolig}  % disable illegal ligatures
\fi
\usepackage{bookmark}
\IfFileExists{xurl.sty}{\usepackage{xurl}}{} % add URL line breaks if available
\urlstyle{same}
\hypersetup{
  pdftitle={제공 수치 방법의 사용 안내},
  pdflang={ko-KR},
  colorlinks=true,
  linkcolor={Maroon},
  filecolor={Maroon},
  citecolor={Blue},
  urlcolor={Blue},
  pdfcreator={LaTeX via pandoc}}

\title{제공 수치 방법의 사용 안내}
\usepackage{etoolbox}
\makeatletter
\providecommand{\subtitle}[1]{% add subtitle to \maketitle
  \apptocmd{\@title}{\par {\large #1 \par}}{}{}
}
\makeatother
\subtitle{함수 호출과 수렴 결과 읽기}
\author{}
\date{}

\begin{document}
\maketitle

{
\hypersetup{linkcolor=}
\setcounter{tocdepth}{3}
\tableofcontents
}
\setstretch{1.12}
\section{제공 방법을 호출하는
순서}\label{uxc81cuxacf5-uxbc29uxbc95uxc744-uxd638uxcd9cuxd558uxb294-uxc21cuxc11c}

이 문서는 함수 사용 시 옆에 두는 짧은 안내다. EOS의 열역학적 표현은 본문
제3·7·9장, 양방향 경계값 문제는 제4·5장, 모드 식별과 수렴은 제6·10장에서
전개한다. 여기의 명령만 실행했다고 그 물리적 설명을 대신한 것은 아니다.

\subsection{한 별과 같은 질량의 두
별}\label{uxd55c-uxbcc4uxacfc-uxac19uxc740-uxc9c8uxb7c9uxc758-uxb450-uxbcc4}

\begin{verbatim}
from tools.stellar_methods import Model, solve_star, solve_pair

star = solve_star("../one_star", alpha=0.0, level="fine")
pair = solve_pair("../one_pair", alpha=0.1, level="fine")
print(pair["delta_f_hz"], pair["delta_p1_hz"])
print(pair["physical_claim_status"])
\end{verbatim}

기본 \texttt{Model}은 동봉한 native Hermite 표, 목표 질량
\(2.644606257493485\times10^{30}\) kg, 양의 표면압력 \(10^{20}\) Pa와
중심 시작점 1 m를 사용한다. 두 모드는 \(\ell=2\)의 \(f,p_1\)이다.
\texttt{solve\_pair}는 두 별을 같은 질량에 맞추고 실제 고유함수를
비교하지만 전체 수렴 실험을 수행하지 않는다. 상태는
\path|pending_additional_checks|다.

\texttt{level}은 \texttt{coarse}, \texttt{medium}, \texttt{fine},
\texttt{audit} 중 하나다. 상대 허용오차와 native 구간 최대 단계는 각각
\((10^{-7},0.2)\), \((10^{-9},0.1)\), \((10^{-11},0.05)\),
\((10^{-13},0.025)\)다. 단계 좌표는 구간별 무차원 좌표이며 반경 m 단위의
격자폭이 아니다. \texttt{method="rk4"}는 별도의 단계법 점검이며 같은
상대오차 제어를 수행한다는 뜻이 아니다.

\begin{verbatim}
custom = Model(target_mass_kg=2.6e30, p_cut_pa=1e20, r0_m=1.0)
pair = solve_pair("../another_mass", model=custom)
\end{verbatim}

이렇게 질량을 바꾼 결과에 기본 실험의 수렴 판정을 자동 부여하지 않는다.
새로운 문제의 branch와 검색 구간, 모드 대응 및 경계 민감도를 다시 정해야
한다. 이 인터페이스가 임의의 EOS·질량에서 근의 존재나 스펙트럼 완전성을
보장하지는 않는다.

\subsection{전체 비교를 실행하고
재개한다}\label{uxc804uxccb4-uxbe44uxad50uxb97c-uxc2e4uxd589uxd558uxace0-uxc7acuxac1cuxd55cuxb2e4}

\begin{verbatim}
from tools.stellar_methods import run_study
result = run_study("../my_study", max_new_runs=2)
print(result["status"], result["pending_runs"])
\end{verbatim}

같은 호출을 반복하면 검증된 완료 사례는 다시 적분하지 않고 필요한 다음
사례로 간다. 한 번에 끝내려면 \texttt{max\_new\_runs}를 생략한다. 이
제한은 한 호출의 작업 수만 바꾸며 물리·수치 허용오차를 낮추지 않는다.

전체 preset은 21개 단일별 계산이다. \texttt{main}은 네 정밀도의 두 별,
\texttt{method}는 별도의 RK4 pair, \texttt{boundary}는
중심·표면·연결점의 세 pair, \texttt{mass}는 두 모델 각각의 질량 기여
추정에 쓰는 네 이웃, \texttt{null}은 무변형 계산이다. 한 그룹만 실행해도
전체 필수 목록은 사라지지 않는다.

\begin{verbatim}
python tools/stellar_methods.py --output ../my_study --group main
python tools/stellar_methods.py --output ../my_study
\end{verbatim}

\texttt{study.json}은 현재 판정과 남은 사례, \texttt{convergence.csv}는
두 절대 주파수와 차이, 각 사례 디렉터리는 입력·수치 출력·실행 상태를
담는다. 모두 끝나기 전에는 \texttt{pending}이다. 최종
\path|numerically_resolved_in_fixed_model|은 정의한 비교에서 수치적
분해가 확인되었다는 뜻이며 엄밀한 오차 상한이나 실제 EOS의 물리적
불확실성 범위를 뜻하지 않는다.

\subsection{초기조건과 해 전체를
확인한다}\label{uxcd08uxae30uxc870uxac74uxacfc-uxd574-uxc804uxccb4uxb97c-uxd655uxc778uxd55cuxb2e4}

\begin{verbatim}
from tools.stellar_methods import Model, solve_pair
m = Model(r0_m=100.0, center_order=2)
pair = solve_pair("../series_start", model=m, level="audit")
\end{verbatim}

\texttt{center\_order=2}는 제3·5장에서 유도한 다음 중심항을 제공 함수로
계산한다. 기본값 0은 1 m의 leading 시작이다. 큰 시작점이 다른 EOS 구간을
건너면 거절한다. 임의의 물리 우변에 보편적으로 적용되는 보정이 아니며,
새 설정의 결과에도 전체 수렴 판정이 자동 부여되지 않는다.

\begin{verbatim}
from tools.variational_check import audit_run
check = audit_run("../my_study/L3_a0.0", "../one_integral.json")
\end{verbatim}

전체 \texttt{run\_study}는 네 paired 수준의 적분 검사를 자동 수행하고
\path|variational_checks.json|에 저장한다. 주파수를 재대입하지 않은
scalar 항등식과 표면항은 본문 제5·6장에 있다. 적분 불일치가 수치 여유를
넘으면 \path|integral_consistency_unresolved|로 남긴다. 원래 주파수를
적분값으로 덮어쓰지는 않는다.

\subsection{실패를 성공 기록과 섞지
않는다}\label{uxc2e4uxd328uxb97c-uxc131uxacf5-uxae30uxb85duxacfc-uxc11euxc9c0-uxc54auxb294uxb2e4}

NaN·무한대, 잘못된 설정, 일치하지 않는 소스·입력·출력은 거절한다. 실패한
계산 폴더는 성공 결과로 사용하지 않는다. 같은 입력을 다시 시도하려면
명시적으로 다음을 사용한다.

\begin{verbatim}
python tools/stellar_methods.py --output ../my_study --retry-failed
\end{verbatim}

원래 시도는 \texttt{failed\_attempts/} 아래에 보존된다. 코드를
수정했거나 물리 입력을 바꿨으면 새 출력 폴더를 선택한다. 동일 출력
폴더의 동시 실행은 잠금으로 거절하며 중단되면 운영체제가 잠금을
해제한다. 이 재개 기능은 POSIX 환경을 사용한다.

\subsection{Fortran에서 가져다
쓴다}\label{fortranuxc5d0uxc11c-uxac00uxc838uxb2e4-uxc4f4uxb2e4}

\begin{verbatim}
use ns_numerical_methods, only: adaptive_integrate, bisect
use ns_stellar_methods, only: integrate_cells, match_cells
\end{verbatim}

첫 모듈은 수치 적분과 근 찾기, 둘째는 EOS 구간에서 TOV·Cowling 물리
함수를 연결하는 adapter다. 실제 인수의 자료형은 해당 모듈과
\texttt{cell\_solve.f90} 예제를 따른다. \texttt{methods}에는 원래 EOS
표를 다루는 완성 함수도 있다. TOV·Cowling 식과 물리 경계조건은
\path|calculation/src/physics/|에 별도로 보인다.

직접 구현한 물리 함수는
\path|make_calculation.py --physics learner/src|로 명시적으로 연결한다.
미완성 함수에 대한 자동 해설 대체는 없다. 제공 예제의 수치적 성공과
학생의 직접 구현 성취도 다른 판정이다.

\end{document}
